\scanpage{63}
% @lines 1–4 / text
aussi, comme $G_0$ opère trivialement sur \uline{tt} module inversible (car $\underline{\operatorname{Hom}}(\mathrm{GP}(1),G_m)=0$) qu'il opère trivialement sur $\underline{\det}(\widetilde{\unsure{\mathfrak g|g}}_0)\simeq\det(E_0)$, donc invarie $\omega$ en plus de $q$.

% @lines 5–9 / text
6) $\mathscr L_0=\mathrm{Gl}(2)$ opérant sur $\mathcal O_{S_0}^2$, opère sur $\mathbb P^1=\mathbb P(\mathcal O_{S_0}^{\unsure{2|?}})$, et comme $G_m$ opère trivialement, on en déduit par passage au quotient une action de $\inserted{G_0=}\mathrm{GP}(1)$ grâce à l'iso (2).

% @lines 10–11 / text
Enfin, l'opération de $G_0$ sur $(E_0,q_0)$ définit une opération sur

% @line 12 / display
\[
(P_0=\check{\mathbb P}(E_0), C_0=\mathbb V(q_0)).
\]
% @below line 12 / text
\noindent\hspace*{.24\linewidth}\tikz[baseline]{\draw (0,0) .. controls (.3,.09) and (.45,-.08) .. (.8,.01) .. controls (1.2,.09) and (1.4,-.04) .. (1.8,.01) .. controls (2.5,.02) and (2.8,-.05) .. (3.2,.04) .. controls (3.5,.08) and (3.8,.0) .. (4.1,.02);}

% @lines 13–14 / text
Nous admettrons \inserted{(pour ne pas nous éterniser...)} que tous les hom obtenus

% @left of lines 13–15 / margin note
\margin{Expliciter le thm d'isomorphie}
% slanted

% @line 15 / display
\[
G_0=\mathrm{GP}(1)\longrightarrow\underline{\operatorname{Aut}}(\unsure{\mathscr X|\mathcal X|\mathfrak X}_0)\tag*{(7)}
\]

% @lines 16–19 / text
sont des \uline{iso}, ce qui en principe fonde le dictionnaire que nous voulons. Mais nous allons \deleted{\illegible{1}} expliciter \inserted{\unsure{de ce dictionnaire|ce dictionnaire|ou dictionnaire}} les points essentiels.

% archive sheet number: 62
